\documentclass[11pt]{article}
\usepackage[T1]{fontenc}
\usepackage[utf8]{inputenc}
\usepackage{lmodern}
\usepackage{amsmath,amssymb,amsthm,mathtools}
\usepackage[margin=1in]{geometry}
\usepackage{booktabs}
\usepackage{enumitem}
\usepackage[colorlinks=true,linkcolor=blue,citecolor=blue,urlcolor=blue]{hyperref}
\usepackage{microtype}
\numberwithin{equation}{section}
\newtheorem{theorem}{Theorem}[section]
\newtheorem{lemma}[theorem]{Lemma}
\newtheorem{proposition}[theorem]{Proposition}
\newtheorem{corollary}[theorem]{Corollary}
\theoremstyle{definition}
\newtheorem{definition}[theorem]{Definition}
\newtheorem{example}[theorem]{Example}
\theoremstyle{remark}
\newtheorem{remark}[theorem]{Remark}
\newcommand{\R}{\mathbb R}
\newcommand{\N}{\mathbb N}
\newcommand{\gph}{\operatorname{gph}}
\newcommand{\dom}{\operatorname{dom}}
\newcommand{\ran}{\operatorname{ran}}
\newcommand{\dist}{\operatorname{dist}}
\newcommand{\diam}{\operatorname{diam}}
\newcommand{\Fix}{\operatorname{Fix}}
\newcommand{\Lip}{\operatorname{Lip}}
\newcommand{\conv}{\operatorname{conv}}
\newcommand{\Id}{\operatorname{Id}}
\newcommand{\rad}{\mathcal R}
\newcommand{\ip}[2]{\langle #1,#2\rangle}
\newcommand{\norm}[1]{\lVert #1\rVert}
\newcommand{\abs}[1]{\lvert #1\rvert}
\newcommand{\eps}{\varepsilon}
\DeclareMathOperator*{\argmin}{argmin}

\title{Sharp Simultaneous Shadows and Fiber Realization\newline for H\"older--RL Relations}
\author{Author information to be supplied}
\date{Research manuscript --- 23 September 2026}

\begin{document}
\maketitle

\begin{abstract}
Let $H$ be a real Hilbert space and let $F:H\rightrightarrows H$ satisfy
$\norm{(x-y)-\lambda(v-w)}\le L\norm{(x-y)+\lambda(v-w)}^\gamma$
for all graph pairs, where $\lambda,L>0$ and $0<\gamma<1$.
We prove that one strongly monotone bi-Lipschitz homeomorphism
$A:H\to H$ simultaneously approximates every nonempty forward fiber
$F(x)$ and inverse fiber $F^{-1}(v)$. Writing
$\rad=L^{1/(1-\gamma)}$, the respective uniform errors are
$\rad/(\sqrt2\lambda)$ and $\rad/\sqrt2$, and the
dimension-independent factor $1/\sqrt2$ is optimal. The construction
also provides a cross inequality and an explicit monotone-enlargement
inclusion. Preserving a specified graph point changes the optimal factor
to $1$. In finite dimensions, for relations that are graph-maximal at the
fixed parameters, the possible complete inverse fibers are exactly the
nonempty compact sets of diameter at most $\rad$, with a scaled statement
for forward fibers. The realization attains the smallest H\"older constant
permitted by the prescribed set. We establish sharp range localization
in fixed windows and give a finite-data quadratic-program construction
of globally invertible surrogates with coverage-dependent paired error
certificates. A deadband model illustrates stable approximate inversion
without claiming stable exact recovery. The results distinguish global
simultaneous organization from the unrestricted compact geometry possible
within an individual fiber.
\end{abstract}

\noindent\textbf{Keywords.} Set-valued relation; H\"older continuity; strong
monotonicity; Kirszbraun extension; fiber realization; certified approximation.

\tableofcontents

\section{Introduction}\label{sec:intro}

A relation may assign several outputs to one input and several inputs to
one output. Describing an individual value set does not necessarily explain
how different value sets fit together. Conversely, the existence of a
well-conditioned approximation need not restrict the topology of the
original fibers. This paper studies these complementary questions for
relations satisfying a sublinear H\"older condition in graph coordinates.

Throughout, $H$ is a real Hilbert space and $\lambda,L>0$ and
$0<\gamma<1$ are fixed. We consider the all-pairs condition
\begin{equation}\label{eq:RLintro}
 \norm{(x-y)-\lambda(v-w)}
 \le L\norm{(x-y)+\lambda(v-w)}^\gamma
 \qquad ((x,v),(y,w)\in\gph F).
\end{equation}
We call this the H\"older--RL condition, where the terminology is only a
label for \eqref{eq:RLintro}. It is a global condition on all graph pairs,
not a condition only relative to a solution or along an individual branch.
The natural length scale is $\rad=L^{1/(1-\gamma)}$.

Our first result concerns approximation at the \emph{same original base
point}. Ordinary graph proximity allows both coordinates of a graph point
to move. In contrast, a forward fiber estimate holds with the input fixed,
and an inverse fiber estimate holds with the output fixed. We construct
one strongly monotone bi-Lipschitz homeomorphism $A:H\to H$ satisfying
\begin{equation}\label{eq:intro_bounds}
 \sup_{v\in F(x)}\norm{v-A(x)}\le\frac{\rad}{\sqrt2\lambda},
 \qquad
 \sup_{x\in F^{-1}(v)}\norm{x-A^{-1}(v)}\le\frac{\rad}{\sqrt2},
\end{equation}
whenever the indicated fiber is nonempty. The same map works for every
input and every output. The factor $1/\sqrt2$ is optimal uniformly over
dimensions, even if the approximating center is not required to come from
a monotone map. We call such an $A$ a \emph{shadow}. This word does not mean
that $A$ is a selection of $F$: its graph need not lie in $\gph F$.

The proof passes through the partial map
$C(x+\lambda v)=x-\lambda v$. A H\"older estimate bounds its excess above
a chosen Lipschitz slope. Orthogonal lifting and the classical
Kirszbraun theorem then give a contraction $N$ with the two-parameter
estimate
\[
 \norm{C(p)-N(q)}^2
 \le \sigma^2\norm{p-q}^2+M_\sigma/2,
 \qquad p\in\dom C,\ q\in H.
\]
It is this cross estimate, rather than just a same-parameter uniform
approximation bound, that yields \eqref{eq:intro_bounds} without a
condition-number loss. Global graph maximality is not needed for this
part of the theory.

Our second main result determines the possible complete fibers of
graph-maximal relations in finite dimensions. At fixed parameters, a set
$K\subset\R^n$ occurs as $F^{-1}(0)$ if and only if it is nonempty,
compact, and has diameter at most $\rad$. The forward-fiber threshold is
$\rad/\lambda$. The sufficient direction constructs a map with prescribed
fixed-point set $K$ and H\"older constant $(\diam K)^{1-\gamma}$, the least
constant compatible with that set. The construction can also make the
original $F$ single-valued and globally Lipschitz. Thus a common stable
shadow is compatible with arbitrary compact fiber geometry at the
permitted scale. The realization theorem is a positive classification,
not merely a collection of examples showing failures of convexity.

We give two further quantitative consequences. First, an exact scalar
localization function gives sharp ball-to-ball range coverage, including
the setting where graph maximality is known only inside a fixed window.
Second, finitely many graph observations produce one globally consistent
contraction through a strongly convex quadratic program. Parameter-space
coverage turns sampled estimates into paired certificates for unobserved
graph points. A finite-information obstruction explains why finitely many
observations cannot certify a uniform error on an unrestricted unbounded
domain. A deadband inverse problem illustrates the precise tradeoff:
stable approximate representatives can coexist with unstable exact inverse
selections. The certificate does not recover information absent from the
original model.

\paragraph{Relation to earlier work.}
Luke and Tam \cite[Theorem~1]{LukeTam} study generalized monotonicity and the proximal
point algorithm, including representations of value sets and inverse value
sets by lower-level sets of lower semicontinuous weakly convex functions
in their relatively maximal submonotone framework. Their original-
relation viewpoint motivates our questions. We consider a different,
sublinear all-pairs condition and ask for simultaneous quantitative
approximation and exact fiber realization. We do not claim to settle the
Bregman, stochastic, or algorithmic directions in their outlook.

The tools underlying our proofs are classical. The Kirszbraun--Valentine
theorem \cite{Kirszbraun,Valentine} supplies Hilbert-space Lipschitz
extension. Same-constant
H\"older extension follows from the Hilbert snowflake embedding associated
with Schoenberg's theorem \cite{Schoenberg}, and is discussed together with
monotone extension by Minty \cite{Minty1970}. Uniform Lipschitz approximation
of uniformly continuous Hilbert-valued mappings is already established
by Levy and Rice \cite{LevyRice}. Constructive links between Lipschitz
extension and monotone operator theory are developed, for example, by
Bauschke and Wang \cite{BauschkeWang}. We do not claim these approximation or
extension principles as new. The conclusion here concerns the original
relation and its inverse at fixed original base points, one common
strongly monotone homeomorphism, and a dimension-independent optimal
constant. The enlargement consequence uses the established notion of
monotone enlargement \cite{BurachikIusemSvaiter}.

The finite-data formula is a specialization of constructive Kirszbraun
extension. Azagra, Le Gruyer, and Mudarra
\cite[Theorems~2 and~8]{AzagraLeGruyerMudarra} give an explicit Lipschitz
extension formula and an extension theorem for strongly bi-Lipschitz maps.
Algorithmic Kirszbraun extension for vector-valued regression is studied
by Zaichyk et al.\ \cite{ZaichykEtAl}. We provide a direct proof of the
particular quadratic program used here, rather than claiming a new
general extension algorithm.

Exact interpolation within prescribed operator classes is a related but
different question. Rubbens, Hendrickx, and Taylor
\cite[Section~3.3]{RubbensHendrickxTaylor} develop strengthened finite-data
conditions for strongly monotone Lipschitz operators, including necessary
relaxations of one-point extendibility. Here the original RL samples need
not be exactly interpolable by that target class. Our construction accepts
a quantified approximation error and provides paired, coverage-dependent
certificates for the same surrogate of the underlying relation. Classical
extension, contraction, and
degree arguments are identified as tools rather than additional novelty
claims.

The present paper is a structure and approximation study. It does not
assert that a shadow together with a residual error bound automatically
implies convergence of an algorithm for the original relation.

\section{Relations, graph coordinates, and fixed parameters}\label{sec:setup}

For $F:H\rightrightarrows H$, write
\[
 \gph F=\{(x,v):v\in F(x)\},\quad
 \dom F=\{x:F(x)\ne\varnothing\},\quad
 \ran F=\{v:F^{-1}(v)\ne\varnothing\},
\]
where $F^{-1}(v)=\{x:v\in F(x)\}$. A \emph{fiber} is one of these
forward or inverse values. An open ball is denoted by $B(x,r)$ and a closed
ball by $\overline B(x,r)$. For nonempty sets $E,D$, their Hausdorff distance
is
\[
 d_H(E,D)=\max\{\sup_{e\in E}\dist(e,D),\sup_{d\in D}\dist(d,E)\},
\]
with the value $+\infty$ allowed. In particular,
$d_H(E,\{a\})=\sup_{e\in E}\norm{e-a}$.

\begin{definition}\label{def:RL}
A nonempty relation $F$ is a $(\lambda,L,\gamma)$ H\"older--RL relation
if \eqref{eq:RLintro} holds. It is \emph{graph-maximal at the fixed
parameters} if no proper superset of its graph satisfies that same
inequality with the same $\lambda,L,\gamma$.
For $U,W\subset H$, a nonempty graph $G\subset U\times W$ is
\emph{relatively maximal} if no proper superset within $U\times W$
satisfies the fixed-parameter inequality.
\end{definition}

This maximality is not maximal monotonicity. A map $A:H\to H$ is
$\mu$-strongly monotone if
$\ip{A(x)-A(y)}{x-y}\ge\mu\norm{x-y}^2$ for all $x,y$. It is
$\ell$-Lipschitz if $\norm{A(x)-A(y)}\le\ell\norm{x-y}$. A bi-Lipschitz
homeomorphism here means a bijection whose map and inverse are Lipschitz.

\begin{lemma}[Graph coordinates]\label{lem:cayley}
Let $F$ satisfy \eqref{eq:RLintro} and put
\begin{equation}\label{eq:cayley}
 D=\{x+\lambda v:(x,v)\in\gph F\},\qquad
 C(x+\lambda v)=x-\lambda v.
\end{equation}
Then $C:D\to H$ is well-defined and $L$-H\"older with exponent $\gamma$.
Conversely, any such map defines an RL graph by
\begin{equation}\label{eq:reconstruct}
 \gph F=\left\{\left(\frac{p+C(p)}2,
                    \frac{p-C(p)}{2\lambda}\right):p\in D\right\}.
\end{equation}
If $D=H$, this graph is graph-maximal at the fixed parameters.
\end{lemma}

\begin{proof}
If two graph points have the same sum $x+\lambda v$, the right-hand side
of \eqref{eq:RLintro} is zero. Their differences $x-\lambda v$ are therefore
equal. Adding and subtracting the equalities identifies the two original
points. The RL inequality is exactly
$\norm{C(p)-C(q)}\le L\norm{p-q}^\gamma$, and solving for $(x,v)$ gives
\eqref{eq:reconstruct}. Conversely, substitution into \eqref{eq:RLintro}
verifies the condition. If $D=H$ and a pair $(x,v)$ is compatible with the
entire graph, compare it with the existing pair at parameter
$p=x+\lambda v$. Their parameter difference is zero, so compatibility
forces their second graph coordinate to agree. The proposed pair is
already in the graph, proving maximality.
\end{proof}

We use two classical extension results. The Kirszbraun--Valentine theorem states that
a Lipschitz map from a subset of a Hilbert space into another Hilbert
space has a same-constant extension to the whole domain space
\cite{Kirszbraun,Valentine}. We also use the following H\"older consequence.

\begin{lemma}[Classical H\"older extension]\label{lem:holderextension}
Every $L$-H\"older map $C:D\subset H\to H$ of exponent $0<\gamma<1$
has an $L$-H\"older extension $\widehat C:H\to H$ of the same exponent.
Consequently an RL relation is graph-maximal if and only if $D=H$ in
\eqref{eq:cayley}. Every RL graph has a graph-maximal completion at the
same parameters.
\end{lemma}

\begin{proof}
The Hilbert snowflake theorem \cite{Schoenberg} gives a Hilbert space $E$
and a map $J:H\to E$ such that
$\norm{J(p)-J(q)}_E=\norm{p-q}^\gamma$.
The map $J(p)\mapsto C(p)$ on $J(D)$ is $L$-Lipschitz, so Kirszbraun
extension to $E$, followed by composition with $J$, gives $\widehat C$.
This proves the extension statement (see also \cite{Minty1970}).
If an RL graph were maximal with $D\ne H$, \eqref{eq:reconstruct} applied
to $\widehat C$ would give a proper compatible extension, a contradiction.
The converse is Lemma~\ref{lem:cayley}, and the same construction proves
the completion assertion.
\end{proof}

\begin{lemma}[Fiber diameters]\label{lem:diameters}
Every RL relation satisfies
\begin{equation}\label{eq:diameters}
 \diam F(x)\le\rad/\lambda,\qquad
 \diam F^{-1}(v)\le\rad
\end{equation}
for each nonempty fiber. Neither maximality nor closedness is needed.
\end{lemma}

\begin{proof}
For $v,w\in F(x)$, set $d=\lambda\norm{v-w}$. Then $d\le Ld^\gamma$.
If $d>0$, division gives $d\le L^{1/(1-\gamma)}=\rad$; the case $d=0$
is immediate. Setting the outputs equal proves the inverse assertion.
\end{proof}

\section{Sharp simultaneous strongly monotone shadows}\label{sec:shadow}

\subsection{Quadratic excess and orthogonal lifting}

For $\sigma>0$, define
\begin{equation}\label{eq:M}
 M_\sigma=(1-\gamma)\gamma^{\gamma/(1-\gamma)}
            \rad^2\sigma^{-2\gamma/(1-\gamma)}.
\end{equation}

\begin{lemma}[Exact quadratic excess]\label{lem:excess}
For every $\sigma>0$,
\begin{equation}\label{eq:excess}
 \sup_{t\ge0}\{L^2t^{2\gamma}-\sigma^2t^2\}=M_\sigma.
\end{equation}
The maximum is attained at
$t_\sigma=(\gamma L^2/\sigma^2)^{1/[2(1-\gamma)]}$.
\end{lemma}

\begin{proof}
For $t>0$ the derivative vanishes precisely when
$\gamma L^2t^{2\gamma-2}=\sigma^2$. The function is positive for
sufficiently small positive $t$, equals zero at zero, and tends to
$-\infty$ at infinity. The unique positive critical point is therefore
the maximum. Since $\sigma^2t_\sigma^2=\gamma L^2t_\sigma^{2\gamma}$,
substitution gives $(1-\gamma)L^2t_\sigma^{2\gamma}$, which simplifies to
\eqref{eq:M}.
\end{proof}

\begin{lemma}[Cross-lift]\label{lem:crosslift}
Let $C:D\subset H\to H$ be $L$-H\"older of exponent $\gamma$ and
$D\ne\varnothing$. For every $0<\sigma<1$ there exists a
$\sigma$-Lipschitz map $N:H\to H$ such that
\begin{equation}\label{eq:crosslift}
 \norm{C(p)-N(q)}^2
 \le\sigma^2\norm{p-q}^2+M_\sigma/2
 \qquad(p\in D,\ q\in H).
\end{equation}
\end{lemma}

\begin{proof}
Let $\ell_2(D)$ be the Hilbert space of square-summable families indexed
by the set $D$, with coordinate unit vectors $e_p$, and set
$r^2=M_\sigma/(2\sigma^2)$. On
$\widehat D=\{(p,re_p):p\in D\}\subset H\oplus\ell_2(D)$ define
$\widehat C(p,re_p)=C(p)$. For distinct $p,s\in D$,
\[
 \norm{C(p)-C(s)}^2
 \le L^2\norm{p-s}^{2\gamma}
 \le\sigma^2\norm{p-s}^2+M_\sigma
 =\sigma^2\norm{(p,re_p)-(s,re_s)}^2.
\]
Thus $\widehat C$ is $\sigma$-Lipschitz; the identical-point case is
automatic. Kirszbraun's theorem gives a $\sigma$-Lipschitz extension
$\widehat N:H\oplus\ell_2(D)\to H$. Set $N(q)=\widehat N(q,0)$.
Comparing a lifted point with a base-plane point gives
\[
 \norm{C(p)-N(q)}^2
 \le\sigma^2(\norm{p-q}^2+r^2)
 =\sigma^2\norm{p-q}^2+M_\sigma/2,
\]
as claimed. Possible nonseparability of $\ell_2(D)$ does not change the
Hilbert-space extension theorem.
\end{proof}

\subsection{The main approximation theorem}

\begin{theorem}[One shadow in both directions]\label{thm:shadow}
Let $F$ be a nonempty $(\lambda,L,\gamma)$ RL relation on a real Hilbert
space. For each $0<\sigma<1$ there exists a bi-Lipschitz homeomorphism
$A_\sigma:H\to H$ satisfying
\begin{align}
 \Lip(A_\sigma)&\le\frac{1+\sigma}{\lambda(1-\sigma)},
 &\Lip(A_\sigma^{-1})&\le\frac{\lambda(1+\sigma)}{1-\sigma},
 \label{eq:conditioning}\\
 \ip{A_\sigma(x)-A_\sigma(y)}{x-y}
 &\ge\frac{1-\sigma}{\lambda(1+\sigma)}\norm{x-y}^2.
 \label{eq:strongmono}
\end{align}
For every $(x,v)\in\gph F$ and $z\in H$, the same map satisfies
\begin{equation}\label{eq:originalcross}
 \ip{x-z}{v-A_\sigma(z)}
 \ge\frac{1-\sigma^2}{4\lambda}
       \norm{x-z+\lambda(v-A_\sigma(z))}^2
       -\frac{M_\sigma}{8\lambda}.
\end{equation}
In particular, with
$r_\sigma=\sqrt{M_\sigma/[2(1-\sigma^2)]}$,
\begin{equation}\label{eq:pairedsigma}
 \norm{v-A_\sigma(x)}\le r_\sigma/\lambda,
 \qquad
 \norm{x-A_\sigma^{-1}(v)}\le r_\sigma
 \qquad((x,v)\in\gph F).
\end{equation}
The displayed radius $r_\sigma$ is minimized at $\sigma=\sqrt\gamma$.
For the resulting map $A$,
\begin{equation}\label{eq:pairedsharp}
 \norm{v-A(x)}\le\frac{\rad}{\sqrt2\lambda},
 \qquad
 \norm{x-A^{-1}(v)}\le\frac{\rad}{\sqrt2}
 \qquad((x,v)\in\gph F).
\end{equation}
\end{theorem}

\begin{proof}
Apply Lemma~\ref{lem:crosslift} directly to the partial map in
Lemma~\ref{lem:cayley}, and define
\begin{equation}\label{eq:shadowparam}
 X(q)=\frac{q+N(q)}2,\qquad
 V(q)=\frac{q-N(q)}{2\lambda}.
\end{equation}
For any $x\in H$, the equation $X(q)=x$ is equivalent to
$q=2x-N(q)$. Its right-hand side is a contraction of modulus at most
$\sigma$, so the Banach fixed-point theorem gives a unique solution.
For any $v\in H$, the same theorem applied to $q=2\lambda v+N(q)$ gives
a unique solution of $V(q)=v$. Thus both $X$ and $V$ are bijections and
$A_\sigma=V\circ X^{-1}$ is a bijection.

For $q_1,q_2\in H$, put $h=q_1-q_2$ and $k=N(q_1)-N(q_2)$, so that
$\norm{k}\le\sigma\norm{h}$. Differences in the original coordinates
are $(h+k)/2$ and $(h-k)/(2\lambda)$. The estimates
\[
 (1-\sigma)\norm{h}\le\norm{h\pm k}
 \le(1+\sigma)\norm{h}
\]
give both bounds in \eqref{eq:conditioning}. Moreover,
\begin{align*}
 \ip{V(q_1)-V(q_2)}{X(q_1)-X(q_2)}
 &=\frac{\norm{h}^2-\norm{k}^2}{4\lambda}\\
 &\ge\frac{1-\sigma^2}{4\lambda}\norm{h}^2
 \ge\frac{1-\sigma}{\lambda(1+\sigma)}
       \norm{X(q_1)-X(q_2)}^2.
\end{align*}
This proves \eqref{eq:strongmono}; the Lipschitz bounds also establish
the asserted homeomorphism property.

For an original graph point and a point on the shadow, let
$p=x+\lambda v\in D$ and $q=z+\lambda A_\sigma(z)$. Then
\[
 p-q=(x-z)+\lambda(v-A_\sigma(z)),\qquad
 C(p)-N(q)=(x-z)-\lambda(v-A_\sigma(z)).
\]
The difference of their squared norms is
$4\lambda\ip{x-z}{v-A_\sigma(z)}$. Therefore
\eqref{eq:crosslift} gives \eqref{eq:originalcross}.
Taking $z=x$ makes its left side zero and gives the forward bound in
\eqref{eq:pairedsigma}. Taking $z=A_\sigma^{-1}(v)$ gives the inverse
bound. These two substitutions use the same $A_\sigma$.

Finally, putting $u=\sigma^2$, the squared radius is a positive
$\sigma$-independent constant times
$u^{-\gamma/(1-\gamma)}(1-u)^{-1}$. Its logarithmic derivative is
\[
 -\frac{\gamma}{(1-\gamma)u}+\frac1{1-u},
\]
which vanishes only at $u=\gamma$. The function diverges at both
endpoints, so this point is the unique minimum. Substitution gives
$M_{\sqrt\gamma}=(1-\gamma)\rad^2$ and $r_{\sqrt\gamma}^2=\rad^2/2$,
proving \eqref{eq:pairedsharp}.
\end{proof}

The quantifiers in this theorem matter: one $A$ controls all graph points.
In an arbitrary Hilbert space the theorem says nothing about the
nonemptiness of a fiber not already represented in $\gph F$.

\begin{corollary}[Enlargement inclusion]\label{cor:enlargement}
The map $A_\sigma$ in Theorem~\ref{thm:shadow} is maximal monotone. With
the standard enlargement
\[
 A_\sigma^{[\eps]}(x)=
 \{v:\ip{x-z}{v-A_\sigma(z)}\ge-\eps\text{ for every }z\in H\},
\]
one has
\begin{equation}\label{eq:enlargement}
 F(x)\subset A_\sigma^{[M_\sigma/(8\lambda)]}(x).
\end{equation}
For the map minimizing the radius in \eqref{eq:pairedsigma},
$\eps=(1-\gamma)\rad^2/(8\lambda)$.
\end{corollary}

\begin{proof}
Strong monotonicity implies monotonicity. If $(u,w)$ is monotonically
related to the entire graph of the continuous full-domain map $A_\sigma$,
then, for $z=u+th$ and $t>0$,
$\ip{h}{w-A_\sigma(u+th)}\le0$. Letting $t\downarrow0$ and then
replacing $h$ by $-h$ gives
$\ip{h}{w-A_\sigma(u)}=0$ for every $h$, whence $w=A_\sigma(u)$.
Thus no point can be added monotonically. Dropping the nonnegative
squared term in \eqref{eq:originalcross} proves the inclusion, and
substituting $\sigma=\sqrt\gamma$ gives the stated parameter.
\end{proof}

The displayed enlargement parameter is associated with the map that
minimizes the fiber-radius guarantee. We do not assert that it is the
smallest possible enlargement parameter.

\subsection{Dimension-independent optimality}

\begin{theorem}[Optimal universal factor]\label{thm:sharpness}
The factor $1/\sqrt2$ in each inequality of \eqref{eq:pairedsharp} cannot
be replaced by a smaller constant valid in all finite dimensions. This
remains true for graph-maximal RL relations and arbitrary choices of a
single center for the fiber under consideration.
\end{theorem}

\begin{proof}
Fix $n\ge1$ and take vertices $u_0,\ldots,u_n\in\R^n$ of a regular
simplex of edge length $\rad$, centered so that $\sum_i u_i=0$.
Let $Q=\conv\{u_0,\ldots,u_n\}$. To compute the vertex radius, use
\[
 \sum_{j=0}^n\norm{u_i-u_j}^2
 =(n+1)\norm{u_i}^2+\sum_{j=0}^n\norm{u_j}^2.
\]
The left side is $n\rad^2$ for every $i$, so all vertex norms are equal
and their common squared value is
$r_n^2=n\rad^2/[2(n+1)]$. For any $b\in\R^n$,
\[
 \frac1{n+1}\sum_{i=0}^n\norm{u_i-b}^2=r_n^2+\norm{b}^2.
\]
At least one vertex is therefore at distance at least $r_n$ from $b$.

The metric projection $P_Q$ is nonexpansive and has image of diameter
$\rad$. Consequently
\[
 \norm{P_Q(p)-P_Q(q)}
 \le\min\{\norm{p-q},\rad\}
 \le\rad^{1-\gamma}\norm{p-q}^\gamma
 =L\norm{p-q}^\gamma.
\]
For the middle inequality, if $t\le\rad$ then
$t\le\rad^{1-\gamma}t^\gamma$, while if $t\ge\rad$ then
$\rad\le\rad^{1-\gamma}t^\gamma$.

Use $C=-P_Q$ in \eqref{eq:reconstruct} with domain $\R^n$. The resulting
relation is graph-maximal by Lemma~\ref{lem:cayley}. Its input is zero
exactly when $p=P_Qp$, or $p\in Q$, and then its output is $p/\lambda$.
Thus $F(0)=Q/\lambda$. Any center $b$ has worst-case distance at least
$r_n/\lambda$ from this fiber. Any universal forward factor $c$ must
therefore satisfy $c\ge\sqrt{n/[2(n+1)]}$ for every $n$, so
$c\ge1/\sqrt2$.

Using $C=P_Q$ instead yields $F^{-1}(0)=Q$ and gives the same lower
bound for any inverse-fiber center. The two witnesses need not be the
same relation: either one already rules out a smaller common universal
factor.
\end{proof}

This theorem does not determine the best factor in each fixed dimension,
nor does it assert optimality of the conditioning bounds in
\eqref{eq:conditioning}--\eqref{eq:strongmono}.

\subsection{Preserving a specified graph point}

\begin{theorem}[Anchored shadow]\label{thm:anchor}
Under the assumptions of Theorem~\ref{thm:shadow}, fix
$(x_0,v_0)\in\gph F$. For every $0<\sigma<1$ there exists a bi-Lipschitz
homeomorphism $A_{\sigma,0}:H\to H$ with the conditioning and
strong-monotonicity bounds
\eqref{eq:conditioning}--\eqref{eq:strongmono}, such that
$A_{\sigma,0}(x_0)=v_0$ and, for every $(x,v)\in\gph F$ and $z\in H$,
\begin{equation}\label{eq:anchorcross}
 \ip{x-z}{v-A_{\sigma,0}(z)}
 \ge\frac{1-\sigma^2}{4\lambda}
      \norm{x-z+\lambda(v-A_{\sigma,0}(z))}^2
      -\frac{M_\sigma}{4\lambda}.
\end{equation}
Writing $A_0=A_{\sqrt\gamma,0}$, this gives, for every original graph point,
\begin{equation}\label{eq:anchorbounds}
 \norm{v-A_{0}(x)}\le\rad/\lambda,
 \qquad \norm{x-A_{0}^{-1}(v)}\le\rad,
\end{equation}
and $F(x)\subset A_0^{[(1-\gamma)\rad^2/(4\lambda)]}(x)$.
The factor $1$ is optimal under the anchoring requirement, already in
dimension one.
\end{theorem}

\begin{proof}
Let $p_0=x_0+\lambda v_0$ and $c_0=x_0-\lambda v_0$. Repeat the lift of
Lemma~\ref{lem:crosslift} with $r^2=M_\sigma/\sigma^2$, and add the base
datum $(p_0,0)\mapsto c_0$. Distinct lifted data are compatible with
Lipschitz modulus $\sigma$ because their extra squared distance supplies
$2M_\sigma$, at least the excess required by Lemma~\ref{lem:excess}.
A lifted datum and the anchor are compatible because
\[
 \norm{C(p)-c_0}^2
 \le\sigma^2\norm{p-p_0}^2+M_\sigma
 =\sigma^2\norm{(p,re_p)-(p_0,0)}^2.
\]
Kirszbraun extension now gives a contraction $N$ with $N(p_0)=c_0$ and
\[
 \norm{C(p)-N(q)}^2\le\sigma^2\norm{p-q}^2+M_\sigma.
\]
Construct $A_{\sigma,0}$ as in \eqref{eq:shadowparam}. The anchor
equalities give $X(p_0)=x_0$ and $V(p_0)=v_0$. The conditioning proof is
unchanged, and the difference-of-squares identity gives
\eqref{eq:anchorcross}. Setting $z=x$ or $z=A_{\sigma,0}^{-1}(v)$ yields
forward and inverse radii
$\lambda^{-1}\sqrt{M_\sigma/(1-\sigma^2)}$ and
$\sqrt{M_\sigma/(1-\sigma^2)}$, respectively. At
$\sigma=\sqrt\gamma$ these equal $\rad/\lambda$ and $\rad$.
Dropping the square gives the stated
enlargement inclusion.

For forward optimality in $\R$, take $Q=[0,\rad]$, use $C=-P_Q$, and
anchor at $(0,0)$. Then $F(0)=[0,\rad/\lambda]$, whereas the anchoring
requires $A_0(0)=0$. Its worst-case forward error is at least
$\rad/\lambda$. For inverse optimality use $C=P_Q$, so
$F^{-1}(0)=Q$, and the same anchor forces $A_0^{-1}(0)=0$.
The worst-case inverse error is at least $\rad$.
\end{proof}

\section{Finite-dimensional geometry and exact fiber realization}
\label{sec:fibers}

\subsection{Full domain, full range, and compact fibers}

\begin{theorem}\label{thm:finite_geometry}
Let $H=\R^n$ and let $F$ be graph-maximal at the fixed RL parameters.
Then
\begin{equation}\label{eq:fullrange}
 \dom F=\ran F=\R^n.
\end{equation}
Every forward and inverse fiber is compact. Both coordinate projections
from $\gph F$ are proper, and $F$ and $F^{-1}$ are upper semicontinuous
and locally bounded. They map compact sets to compact unions of values.
Orient the graph by the parametrization of Lemma~\ref{lem:cayley},
with the standard orientation on $\R^n$. Under this convention,
each projection has degree $1$.
\end{theorem}

\begin{proof}
By Lemma~\ref{lem:holderextension}, the Cayley map $C$ is defined on all
of $\R^n$. It is continuous and satisfies
$\norm{C(p)}\le\norm{C(0)}+L\norm{p}^\gamma$.
The map
$p\mapsto((p+C(p))/2,(p-C(p))/(2\lambda))$
is a homeomorphism onto the graph, with continuous inverse $(x,v)\mapsto
x+\lambda v$. For $t\in[0,1]$ define
\[
 X_t(p)=\frac{p+tC(p)}2,\qquad
 V_t(p)=\frac{p-tC(p)}{2\lambda}.
\]
Uniformly in $t$,
\[
 2\norm{X_t(p)}\ge\norm{p}-\norm{C(0)}-L\norm{p}^\gamma,
 \quad
 2\lambda\norm{V_t(p)}\ge\norm{p}-\norm{C(0)}-L\norm{p}^\gamma.
\]
The right sides tend to infinity with $\norm p$. Hence both homotopies
are uniformly proper: preimages of any bounded target set remain in a
common bounded parameter set.

Fix a target $y$. Choose a large ball so that neither homotopy hits $y$
on its boundary. Homotopy invariance of Brouwer degree gives degree $1$
for $X_1$ and $V_1$, since $X_0=\Id/2$ and $V_0=\Id/(2\lambda)$ are
positive scalar multiples of the identity. The nonzero-degree existence
property gives a preimage of $y$ for each map, proving \eqref{eq:fullrange}.
These are classical degree properties; see, for example,
\cite{BenevieriFuriPeraSpadini}.

The displayed coercive bounds and continuity show that the preimage of
a compact set under either projection, expressed in the parameter $p$,
is closed and bounded, hence compact. The projections are thus proper,
and point preimages give compact fibers. The graph is closed: if
$(x_j,v_j)\to(x,v)$, then $p_j=x_j+\lambda v_j\to x+\lambda v$ and
$C(p_j)=x_j-\lambda v_j\to x-\lambda v$, which verifies graph membership
by continuity. The same coercive bounds show local boundedness of each
fiber map: bounded $x$ bounds $p$, hence $v$, and bounded $v$ bounds $p$,
hence $x$.

For upper semicontinuity, suppose an open set $O$ contains $F(x)$ but
there are $x_j\to x$ and $v_j\in F(x_j)\setminus O$. Local boundedness
gives a convergent subsequence of $v_j$; closedness of the graph puts its
limit in $F(x)\subset O$, contradicting openness. The inverse argument
is identical with the coordinates exchanged. Finally, for a compact
set $K$, the portion of the graph above $K$ is compact by properness;
its projection onto the other coordinate is $F(K)$ and is compact.
The inverse statement follows from the other projection.
\end{proof}

\begin{corollary}\label{cor:hausdorff}
For a relation satisfying Theorem~\ref{thm:finite_geometry}, the same
shadow $A$ satisfies the complete-fiber Hausdorff estimates
\[
 \sup_{x\in\R^n}d_H(F(x),\{A(x)\})\le\frac{\rad}{\sqrt2\lambda},
 \qquad
 \sup_{v\in\R^n}d_H(F^{-1}(v),\{A^{-1}(v)\})\le\frac{\rad}{\sqrt2}.
\]
\end{corollary}
\begin{proof}
All fibers are nonempty by \eqref{eq:fullrange}, and Hausdorff distance
to a singleton equals the supremum of the distances to that singleton.
Apply \eqref{eq:pairedsharp}.
\end{proof}

Finite dimensionality in Theorem~\ref{thm:finite_geometry} is essential
to this proof: bounded closed subsets of an infinite-dimensional Hilbert
space need not be compact, and Brouwer degree is not available in that
generality. We do not assert \eqref{eq:fullrange} for arbitrary Hilbert
spaces.

\subsection{Prescribed fixed-point sets with optimal H\"older modulus}

\begin{lemma}[A strict convex-hull margin]\label{lem:margin}
Let $K$ be a nonempty compact subset of a Hilbert space, let
$D_K=\diam K$, and let $Q=\overline{\conv}K$. Then, for every $p\in Q$,
\begin{equation}\label{eq:margin}
 \sup_{q\in Q}\norm{p-q}^2
 \le D_K^2-\dist(p,K)^2.
\end{equation}
\end{lemma}

\begin{proof}
Every $k\in K$ and $q\in Q$ satisfy $\norm{k-q}\le D_K$: this follows
by convexity of the norm for a finite convex combination $q$, and then
by passage to its limit. Choose finite convex combinations
$p_m=\sum_i\alpha_i k_i\to p$, with $\alpha_i\ge0$ and
$\sum_i\alpha_i=1$. For any fixed $q\in Q$, the Hilbert variance identity
gives
\[
 \norm{p_m-q}^2
 =\sum_i\alpha_i\norm{k_i-q}^2
  -\sum_i\alpha_i\norm{k_i-p_m}^2
 \le D_K^2-\dist(p_m,K)^2.
\]
The distance function is continuous, so taking $m\to\infty$ proves
the inequality for each $q$. Taking its supremum gives
\eqref{eq:margin}.
\end{proof}

\begin{theorem}[Optimal prescribed fixed-point realization]\label{thm:fixedset}
Let $K$ be a nonempty compact subset of a real Hilbert space and
$D_K=\diam K$. For every $0<\gamma<1$ there exists a map
$R:H\to\overline{\conv}K$ such that
\begin{equation}\label{eq:fixedset}
 \Fix R=K,\qquad
 \norm{R(x)-R(y)}\le D_K^{1-\gamma}\norm{x-y}^\gamma
 \quad(x,y\in H).
\end{equation}
For $D_K=0$, the coefficient is interpreted as zero. For $D_K>0$, no
smaller H\"older constant is possible for any map fixing every point of
$K$.
\end{theorem}

\begin{proof}
If $D_K=0$, the constant map to the sole point of $K$ has all asserted
properties. Assume $D_K>0$, set $Q=\overline{\conv}K$ and
$L_0=D_K^{1-\gamma}$, and fix $k_0\in K$.
If $Q=K$, use $R=P_Q$. Nonexpansiveness and the range diameter give
\[
 \norm{P_Qx-P_Qy}\le\min\{\norm{x-y},D_K\}
 \le L_0\norm{x-y}^\gamma,
\]
and the fixed-point set is $Q=K$.

Suppose $Q\setminus K\ne\varnothing$. For each $p\in Q\setminus K$,
put $d_p=\dist(p,K)>0$ and $E_p=\sqrt{D_K^2-d_p^2}<D_K$. By
Lemma~\ref{lem:margin}, $\norm{p-q}\le E_p$ for every $q\in Q$.
Choose
\[
 0<r_p<\min\{d_p/2,(D_K-E_p)/2\},\qquad D_p=E_p+r_p<D_K,
\]
and define on $Q$
\[
 \phi_p(y)=\max\{1-\norm{y-p}/r_p,0\},\qquad
 H_p(y)=\phi_p(y)(y-k_0),\qquad M_p=1+D_K/r_p.
\]
The function $\phi_p$ vanishes on $K$, is positive in the relative open
ball $Q\cap B(p,r_p)$, and has Lipschitz constant at most $1/r_p$.
For $y,z\in Q$,
\[
 \norm{H_p(y)-H_p(z)}
 \le\phi_p(y)\norm{y-z}
       +\abs{\phi_p(y)-\phi_p(z)}\norm{z-k_0}
 \le M_p\norm{y-z}.
\]
Choose a positive number
\begin{equation}\label{eq:epsilon_bump}
 0<\eps_p\le
 \min\left\{1,\frac{L_0D_p^{\gamma-1}-1}{M_p}\right\}.
\end{equation}
The right side is positive because $D_p<D_K$ and $\gamma-1<0$.
Set $S_p(y)=y-\eps_p H_p(y)$. It is a convex combination of $y$ and
$k_0$, so $S_p(Q)\subset Q$, and it fixes $K$.

We check its H\"older modulus for all point pairs. If
$\phi_p(y)=\phi_p(z)=0$, then $S_p(y)-S_p(z)=y-z$ and
$\norm{y-z}\le D_K$ gives the required $L_0$-H\"older bound. Otherwise
assume, after interchanging $y,z$ if necessary, that $\phi_p(y)>0$.
Then $h=\norm{y-z}<r_p+E_p=D_p$. If $h=0$ there is nothing to prove;
if $h>0$, \eqref{eq:epsilon_bump} gives
\[
 \norm{S_p(y)-S_p(z)}
 \le(1+\eps_p M_p)h
 \le L_0D_p^{\gamma-1}h
 \le L_0h^\gamma.
\]
The last inequality uses the negative exponent $\gamma-1$.

The compact metric space $K$ is separable, and rational convex
combinations of a countable dense subset make $Q$ separable. Hence
$Q\setminus K$ is second countable. Its cover by the relative open
balls $Q\cap B(p,r_p)$ has a finite or countable subcover. Enumerate the
associated maps as $S_j$, and choose strictly positive weights $a_j$
summing to one. Define
\[
 S(y)=\sum_j a_jS_j(y).
\]
The series converges uniformly because $Q$ is bounded. Its sum lies in
$Q$: finite weighted sums completed with the remaining weight at $k_0$
lie in $Q$ and converge to the same limit. The common H\"older estimate
passes to the sum. Moreover,
\[
 S(y)=y-\theta(y)(y-k_0),\qquad
 \theta(y)=\sum_j a_j\eps_j\phi_j(y).
\]
For $y\in K$, every bump vanishes, so $S(y)=y$. If $y\in Q\setminus K$,
one of the covering bumps is positive; all coefficients are positive,
so $\theta(y)>0$. Since $y\ne k_0$, this implies $S(y)\ne y$.
Thus $\Fix S=K$ on $Q$.

Finally set $R=S\circ P_Q$. Nonexpansiveness of $P_Q$ transfers the
$L_0$-H\"older estimate to all of $H$. If $R(x)=x$, then $x\in Q$,
so $P_Qx=x$ and $S(x)=x$, hence $x\in K$. Conversely every point of $K$
is fixed. This proves \eqref{eq:fixedset}.

If any map fixes $K$ and has H\"older constant $B$, then for distinct
$k,l\in K$ one has $\norm{k-l}\le B\norm{k-l}^\gamma$, so
$B\ge\norm{k-l}^{1-\gamma}$. Taking the supremum over $K$ gives
$B\ge D_K^{1-\gamma}$ and proves optimality.
\end{proof}

\begin{theorem}[Exact finite-dimensional fiber classification]\label{thm:fibers}
Fix $n\ge1$ and the parameters $\lambda,L>0$, $0<\gamma<1$.
A set $K\subset\R^n$ is the complete inverse fiber $F^{-1}(0)$ of
some graph-maximal RL relation if and only if
\begin{equation}\label{eq:fiberiff}
 K\ne\varnothing,\qquad K\text{ is compact},\qquad \diam K\le\rad.
\end{equation}
A set is the complete forward fiber $F(0)$ of such a relation if and
only if it is nonempty, compact, and has diameter at most
$\rad/\lambda$. Translations give the corresponding statements for
any specified base input or output.
\end{theorem}

\begin{proof}
Necessity follows from Theorem~\ref{thm:finite_geometry} and
Lemma~\ref{lem:diameters}. For sufficiency in the inverse case, take
the map $R$ from Theorem~\ref{thm:fixedset}. Its H\"older constant is
$(\diam K)^{1-\gamma}\le L$. Use $C=R$ with full domain in
\eqref{eq:reconstruct}. Lemma~\ref{lem:cayley} gives an RL relation
maximal at the specified parameters, even if its least possible
H\"older constant is smaller than $L$. Its output is zero exactly when
$p=R(p)$, and at such a parameter its input equals $p$. Therefore
$F^{-1}(0)=\Fix R=K$.

For a forward fiber $K$ with diameter at most $\rad/\lambda$, apply
Theorem~\ref{thm:fixedset} to $\lambda K$ to obtain $S$ with
$\Fix S=\lambda K$ and constant at most $L$. Use $C=-S$ in
\eqref{eq:reconstruct}. The input is zero exactly when $p=S(p)$,
and then the output is $p/\lambda$. Hence $F(0)=K$. Translating either
original coordinate does not change graph differences in the RL
inequality, proving the last assertion.
\end{proof}

\subsection{A single-valued Lipschitz realization}

\begin{corollary}\label{cor:lipschitz_realization}
Under \eqref{eq:fiberiff}, for every $0<\eta<1$ the inverse fiber can
be realized by a single-valued globally Lipschitz map
$F:\R^n\to\R^n$ that is graph-maximal at the same parameters and
satisfies $F^{-1}(0)=K$. Writing $Q=\overline{\conv}K$ and
$D_K=\diam K$, it may be chosen so that
\begin{equation}\label{eq:lipschitz_realization}
 \Lip F\le\frac{1+\eta}{\lambda(1-\eta)},\qquad
 \sup_u\left\|F(u)-\frac{u-P_Qu}{\lambda}\right\|
 \le\frac{\eta D_K}{\lambda}.
\end{equation}
The least H\"older constant of its Cayley map remains
$D_K^{1-\gamma}$ (zero for a singleton).
\end{corollary}

\begin{proof}
For $D_K=0$, with $K=\{k_0\}$ take $F(u)=(u-k_0)/\lambda$.
All assertions follow directly. For $D_K>0$, in
\eqref{eq:epsilon_bump} impose in addition $\eps_p\le\eta/M_p$.
This preserves all previous conclusions. With $E=S-\Id_Q$ the same
positive combination gives
$\Lip E\le\eta$ and $\sup_{y\in Q}\norm{E(y)}\le\eta D_K$.
The projection case $Q=K$ is included by taking $E=0$.
Thus the realized map can be written
$R=P_Q+E\circ P_Q$.

Define $G=(\Id+R)/2$. Projection onto a closed convex set is monotone
and nonexpansive. Consequently
\[
 2\ip{G(a)-G(b)}{a-b}
 =\norm{a-b}^2+\ip{P_Qa-P_Qb}{a-b}
   +\ip{E(P_Qa)-E(P_Qb)}{a-b}
 \ge(1-\eta)\norm{a-b}^2.
\]
Also $\Lip G\le(2+\eta)/2$. Put $c=(1-\eta)/2$ and
$M=(2+\eta)/2$. For any target $u$ and any $0<t<2c/M^2$, the map
$a\mapsto a-t(G(a)-u)$ is a contraction, since the square of its
Lipschitz modulus is at most $1-2tc+t^2M^2<1$. The fixed-point theorem
therefore shows that $G$ is bijective. Strong monotonicity and
Cauchy--Schwarz give $\Lip G^{-1}\le2/(1-\eta)$.

The relation generated by $R$ is thus the single-valued map
\[
 F(u)=\frac{a-R(a)}{2\lambda}=\frac{a-u}{\lambda},
 \qquad a=G^{-1}(u).
\]
Since $\Id-P_Q$ is nonexpansive,
\[
 \norm{(a-R(a))-(b-R(b))}\le(1+\eta)\norm{a-b},
\]
which together with the inverse bound gives the first part of
\eqref{eq:lipschitz_realization}. Maximality and the exact zero set
were proved in Theorem~\ref{thm:fibers}. Since $R$ fixes $K$, the
optimality argument of Theorem~\ref{thm:fixedset} gives its least
H\"older constant.

It remains to verify the uniform approximation. Set
$G_0=(\Id+P_Q)/2$. Its strong-monotonicity modulus is at least $1/2$,
and its inverse is $G_0^{-1}(u)=2u-P_Qu$. Indeed the projection
characterization says $\ip{u-P_Qu}{y-P_Qu}\le0$ for every $y\in Q$;
multiplying by $2$ shows $P_Q(2u-P_Qu)=P_Qu$.
For $a=G^{-1}(u)$ and $a_0=G_0^{-1}(u)$,
\[
 \norm{a-a_0}
 \le2\norm{G_0(a)-G_0(a_0)}
 =\norm{E(P_Qa)}\le\eta D_K.
\]
Since $(u-P_Qu)/\lambda=(a_0-u)/\lambda$, subtraction gives the
second inequality of \eqref{eq:lipschitz_realization}.
\end{proof}

In particular, nonconvex compact zero sets can occur in a single-valued
globally Lipschitz realization arbitrarily close to the monotone
projection residual $(\Id-P_Q)/\lambda$. This does not make that
realization strongly monotone: a strongly monotone map has at most one
zero.

\section{Sharp range localization in global and fixed windows}
\label{sec:range}

\begin{lemma}[Scalar localization]\label{lem:rho}
For $t\ge0$, let $\rho(t)$ be the largest nonnegative solution of
\begin{equation}\label{eq:rho}
 \rho-t=L(\rho+t)^\gamma.
\end{equation}
Then $\rho(0)=\rad$, $\rho$ is continuous and strictly increasing,
and its range is $[\rad,\infty)$. For every graph anchor
$(x_0,v_0)\in\gph F$, every graph point satisfies
\begin{equation}\label{eq:localization}
 \norm{x-x_0}\le\rho(\lambda\norm{v-v_0}),\qquad
 \lambda\norm{v-v_0}\le\rho(\norm{x-x_0}).
\end{equation}
\end{lemma}

\begin{proof}
Write $u=\rho+t$. Equation \eqref{eq:rho} is
$u-Lu^\gamma=2t$. On $u\ge\rad$, the left side has derivative
$1-\gamma Lu^{\gamma-1}\ge1-\gamma>0$, starts at zero, and tends to
infinity. It thus defines a unique continuous branch $u(t)\ge\rad$.
Its derivative is
$u'(t)=2/(1-\gamma Lu(t)^{\gamma-1})>2$, so
$\rho(t)=u(t)-t$ is strictly increasing and unbounded. For $t>0$ a
nonnegative solution of \eqref{eq:rho} must have $\rho>t$ and
$u-Lu^\gamma>0$, hence $u>\rad$, so this solution is unique. For
$t=0$ the possible roots are $0$ and $\rad$, and the stated choice is
the larger root.

For $s=\norm{x-x_0}$ and $t=\lambda\norm{v-v_0}$, the reverse and
ordinary triangle inequalities give
$\abs{s-t}\le L(s+t)^\gamma$. If $s\le t$, then $s\le\rho(t)$
because $\rho(t)\ge t$. If $s>t$, the inequality says
$(s+t)-L(s+t)^\gamma\le2t$. When $s+t\le\rad$ the desired bound
is automatic from $\rho(t)\ge\rad$; when $s+t>\rad$, monotonicity
of the preceding scalar function gives $s+t\le u(t)$, hence
$s\le\rho(t)$. Interchanging $s,t$ proves the other bound.
\end{proof}

For $r>\rad$ define $h(r)\in(0,r)$ by
\begin{equation}\label{eq:h}
 r-h(r)=L(r+h(r))^\gamma.
\end{equation}
The left side minus the right side is strictly decreasing as a function
of $h$, positive at $h=0$ and negative at $h=r$. Thus the root is
unique. Lemma~\ref{lem:rho} shows $\rho(h(r))=r$.

\begin{theorem}[Global coverage and entire-fiber containment]\label{thm:coverage}
Let $F$ be graph-maximal on $\R^n$, and fix
$(x_0,v_0)\in\gph F$. For $r>\rad$,
\begin{equation}\label{eq:coverage}
 B(v_0,h(r)/\lambda)\subset F(B(x_0,r)),
 \qquad
 F^{-1}(v)\subset B(x_0,r)
 \quad\text{for every }v\in B(v_0,h(r)/\lambda).
\end{equation}
In the second assertion the inverse fiber is nonempty. The symmetric
statement is
\[
 B(x_0,h(\lambda s))\subset F^{-1}(B(v_0,s)),
 \quad F(x)\subset B(v_0,s)
 \quad(x\in B(x_0,h(\lambda s))),
\]
whenever $\lambda s>\rad$.
\end{theorem}

\begin{proof}
Theorem~\ref{thm:finite_geometry} guarantees a nonempty inverse fiber
for every $v$. If $\lambda\norm{v-v_0}<h(r)$, then every point in
that fiber satisfies
\[
 \norm{x-x_0}\le\rho(\lambda\norm{v-v_0})
 <\rho(h(r))=r.
\]
This gives both coverage and full-fiber containment. For the symmetric
statement, full domain gives $F(x)\ne\varnothing$ and the second
inequality of \eqref{eq:localization} gives
$\lambda\norm{v-v_0}<\rho(h(\lambda s))=\lambda s$ for every
$v\in F(x)$.
\end{proof}

\begin{theorem}[Fixed-window coverage]\label{thm:window}
Let $U,W\subset\R^n$ and let $G\subset U\times W$ be a nonempty
relatively maximal RL graph. Suppose $(x_0,v_0)\in G$,
$B(x_0,r)\subset U$, $B(v_0,s)\subset W$, and $r>\rad$, $s>0$.
Then every
\begin{equation}\label{eq:windowtargets}
 v\in B\bigl(v_0,\min\{s,h(r)/\lambda\}\bigr)
\end{equation}
has a nonempty inverse fiber in $G$, and all its inverse points in
$G$ belong to $B(x_0,r)$. In particular, the target ball in
\eqref{eq:windowtargets} is contained in the range of the restriction
of $G$ to $B(x_0,r)\times W$.
\end{theorem}

\begin{proof}
Apply Lemma~\ref{lem:holderextension} to the partial Cayley map of $G$
and let $\widehat F$ be the resulting global graph-maximal completion.
The graph $\gph\widehat F\cap(U\times W)$ contains $G$ and is
compatible at the same parameters. Relative maximality therefore gives
\[
 G=\gph\widehat F\cap(U\times W).
\]
For a target in \eqref{eq:windowtargets},
Theorem~\ref{thm:coverage} gives a nonempty complete inverse fiber of
$\widehat F$, entirely in $B(x_0,r)\subset U$. The target belongs
to $W$, so the equality of graph restrictions puts every one of these
points in $G$. Conversely, every point of $G$ is in the completion,
and is localized by the same theorem. This proves all assertions.
\end{proof}

\begin{proposition}[Sharp radius and threshold]\label{prop:coveragesharp}
The radius $h(r)/\lambda$ in \eqref{eq:coverage} cannot be increased
uniformly, and the strict threshold $r>\rad$ cannot be replaced by
$r\le\rad$ while retaining a positive centered output ball uniformly
over the class.
\end{proposition}

\begin{proof}
In dimension one, take $C(p)=L(p_+)^\gamma$, where
$p_+=\max\{p,0\}$. The inequality
$\abs{a^\gamma-b^\gamma}\le\abs{a-b}^\gamma$ for $a,b\ge0$,
together with nonexpansiveness of $p\mapsto p_+$, makes $C$
$L$-H\"older. Its full graph produces a maximal RL relation anchored
at $(0,0)$, with
\[
 x(p)=\frac{p+L(p_+)^\gamma}{2},\qquad
 v(p)=\frac{p-L(p_+)^\gamma}{2\lambda}.
\]
A positive output requires $p>\rad$, since for $p\le0$ the output
is nonpositive and for $0<p\le\rad$ one has $p\le Lp^\gamma$.
On $p>\rad$, $v$ is strictly increasing because
$v'(p)=(1-\gamma Lp^{\gamma-1})/(2\lambda)>0$, and it ranges over
all positive outputs. The corresponding $x(p)$ is strictly greater
than $\rad$. Hence no positive output has an inverse in
$(-r,r)$ when $r\le\rad$.

For $r>\rad$, the parameter $p=r+h(r)$ satisfies
$L p^\gamma=r-h(r)$, so $x(p)=r$ and
$v(p)=h(r)/\lambda$. This positive output has no other inverse,
and its only inverse is outside the open input ball. Any strictly
larger centered output radius would include this target and fail.
The global example is also a permitted fixed-window example by taking
the windows to be the full space.
\end{proof}

\section{Finite-data shadows and coverage certificates}\label{sec:finite}

The full-graph theorem is an existence theorem: an arbitrary extension
does not by itself supply a finite numerical evaluation oracle. We now
give a separate construction from a fixed finite data set in $\R^n$.
The resulting map is defined globally, but its accuracy for an unknown
original relation is certified only at observations and where coverage
is available.

\subsection{A globally consistent quadratic-program construction}

Fix $m\ge1$ graph samples $(x_i,v_i)$ and set
$p_i=x_i+\lambda v_i$, $c_i=x_i-\lambda v_i$. Repeated parameter
points in consistent data have the same $c_i$ and may be merged. Fix
$0<\sigma<1$, and put
\begin{equation}\label{eq:finiteconstants}
 M=M_\sigma,\qquad a^2=M/2,
 \qquad P=[p_1\ \cdots\ p_m],\quad V=[c_1\ \cdots\ c_m].
\end{equation}
The construction requires only the sample compatibility inequalities
\begin{equation}\label{eq:samplecompat}
 \norm{c_i-c_j}^2\le\sigma^2\norm{p_i-p_j}^2+M.
\end{equation}
Certified H\"older data satisfy these by Lemma~\ref{lem:excess}. Define
\begin{align}
 Q&=\sigma^2P^TP+V^TV+a^2I_m,\label{eq:Q}\\
 d_i(q)&=\sigma^2\norm{p_i}^2-\norm{c_i}^2
              -4\sigma^2\ip{p_i}{q},\label{eq:d}\\
 \Delta_m&=\{\theta\in\R^m:\theta_i\ge0,\ \sum_i\theta_i=1\},\notag\\
 \theta(q)&=\argmin_{\theta\in\Delta_m}
           \Phi_q(\theta),\qquad
 \Phi_q(\theta)=\theta^TQ\theta+d(q)^T\theta,
 \qquad N_m(q)=V\theta(q).\label{eq:QP}
\end{align}

\begin{theorem}[Finite-data contraction and sampled cross estimate]
\label{thm:finite_qp}
The minimizer in \eqref{eq:QP} exists and is unique. The formula defines
one global $\sigma$-Lipschitz map $N_m:\R^n\to\R^n$. Under
\eqref{eq:samplecompat}, it satisfies
\begin{equation}\label{eq:samplecross}
 \norm{c_i-N_m(q)}^2
 \le\sigma^2\norm{p_i-q}^2+a^2
 \qquad(i=1,\ldots,m,\ q\in\R^n).
\end{equation}
The graph parametrization
\begin{equation}\label{eq:Am}
 x=\frac{q+N_m(q)}2,\qquad
 A_m(x)=\frac{q-N_m(q)}{2\lambda}
\end{equation}
defines a strongly monotone bi-Lipschitz homeomorphism with the bounds
\eqref{eq:conditioning}--\eqref{eq:strongmono}. It obeys the cross
inequality \eqref{eq:originalcross} for every sampled graph point and
all $z$, and therefore the paired sampled bounds
\begin{equation}\label{eq:samplepaired}
 \norm{v_i-A_m(x_i)}\le\frac{a}{\lambda\sqrt{1-\sigma^2}},
 \qquad
 \norm{x_i-A_m^{-1}(v_i)}\le\frac{a}{\sqrt{1-\sigma^2}}.
\end{equation}
\end{theorem}

\begin{proof}
The simplex is compact and nonempty, and $Q\succeq a^2I_m$ with
$a>0$. Thus the continuous objective is strictly convex and has a
unique minimizer.

Let $r=a/\sigma$, introduce lifted points $z_i=(p_i,re_i)$ in
$E=\R^n\oplus\R^m$, and let $Z$ have these columns. Then
$Q=\sigma^2 Z^TZ+V^TV$. For every $z\in E$, consider the same
minimization with coefficients
\[
 d_i^E(z)=\sigma^2\norm{z_i}^2-\norm{c_i}^2
                        -4\sigma^2\ip{z_i}{z},
\]
and denote its unique minimizer by $\theta_z$. At $z=(q,0)$ the
coefficients differ from \eqref{eq:d} by the same constant $a^2$
for every $i$. This changes the simplex objective by a constant,
so $\theta_{(q,0)}=\theta(q)$.

For $z,w\in E$, put $\delta=\theta_z-\theta_w$. The variational
inequalities at the two minimizers, tested against one another, give
\[
 \delta^TQ\delta\le2\sigma^2\ip{Z\delta}{z-w}.
\]
Indeed subtraction of the gradients uses
$d^E(z)-d^E(w)=-4\sigma^2Z^T(z-w)$ and division by $2$ gives the
displayed inequality. Hence
\begin{align*}
 \norm{V\delta}^2
 &\le\sigma^2\bigl(2\ip{Z\delta}{z-w}-\norm{Z\delta}^2\bigr)\\
 &=\sigma^2\bigl(\norm{z-w}^2-\norm{Z\delta-(z-w)}^2\bigr)
 \le\sigma^2\norm{z-w}^2.
\end{align*}
Thus $\widetilde N(z)=V\theta_z$ is $\sigma$-Lipschitz, and its
restriction $N_m(q)=\widetilde N(q,0)$ has the same modulus. This
Lipschitz argument did not require sample compatibility.

Compatibility ensures interpolation at the lifted data. At $z=z_j$,
the candidate $\theta=e_j$ satisfies, for each $i$,
\[
 [2Qe_j+d^E(z_j)]_i-[2Qe_j+d^E(z_j)]_j
 =\sigma^2\norm{z_i-z_j}^2-\norm{c_i-c_j}^2\ge0.
\]
For $i\ne j$, the last inequality is \eqref{eq:samplecompat}, since
$\sigma^2\norm{z_i-z_j}^2=\sigma^2\norm{p_i-p_j}^2+2a^2$.
For $i=j$ it is equality. These inequalities give the simplex
first-order optimality condition at $e_j$, so uniqueness implies
$\theta_{z_j}=e_j$ and $\widetilde N(z_j)=c_j$. Comparing $z_j$
with $(q,0)$ yields \eqref{eq:samplecross}.

The contraction-to-operator proof of Theorem~\ref{thm:shadow}, applied
to $N_m$, proves the assertions about $A_m$. The same
difference-of-squares calculation applied to \eqref{eq:samplecross}
gives \eqref{eq:originalcross} on sampled graph points, and setting
the original input or output equal gives \eqref{eq:samplepaired}.
\end{proof}

The quadratic program is also obtained by specializing the explicit
Kirszbraun formula of \cite{AzagraLeGruyerMudarra} to the finitely many
lifted data. The preceding proof is self-contained for this formula and
does not require evaluating a convex envelope. The sample set is fixed:
adding observations may change the resulting global map. Independently
choosing a feasible ball-intersection value at each query, without
cross-query compatibility, is not the construction in
\eqref{eq:QP}.

\subsection{Coverage-dependent paired certificates}

\begin{theorem}\label{thm:covered}
Suppose the samples belong to an RL graph with partial Cayley map $C$.
For $\delta\ge0$ put $b=L\delta^\gamma$ and
\begin{equation}\label{eq:Kdelta}
 K_\delta=
 \frac{b+\sigma^2\delta+
 \sqrt{\sigma^2(b+\delta)^2+(1-\sigma^2)a^2}}
 {1-\sigma^2}.
\end{equation}
Every original graph point $(x,v)$ whose parameter
$p=x+\lambda v$ satisfies
$\dist(p,\{p_1,\ldots,p_m\})\le\delta$ obeys
\begin{equation}\label{eq:coveredpaired}
 \lambda\norm{v-A_m(x)}\le K_\delta,
 \qquad \norm{x-A_m^{-1}(v)}\le K_\delta.
\end{equation}
At $\delta=0$, these are the bounds in
\eqref{eq:samplepaired}; if also $\sigma=\sqrt\gamma$,
$K_0=\rad/\sqrt2$.
\end{theorem}

\begin{proof}
Choose a closest sample $p_i$, which exists because the set is finite.
For every $q\in\R^n$, H\"older continuity and
\eqref{eq:samplecross} give
\begin{equation}\label{eq:coveragecross}
 \norm{C(p)-N_m(q)}
 \le b+\sqrt{\sigma^2(\norm{p-q}+\delta)^2+a^2}.
\end{equation}
First take $q=x+\lambda A_m(x)$, so the compared points have the
same original input. Then
$\norm{p-q}=\norm{C(p)-N_m(q)}=\lambda\norm{v-A_m(x)}$.
For the inverse comparison choose
$q=A_m^{-1}(v)+\lambda v$; both norms instead equal
$\norm{x-A_m^{-1}(v)}$. Thus in either case the nonnegative quantity
$t$ to be bounded satisfies
\[
 t\le b+\sqrt{\sigma^2(t+\delta)^2+a^2}.
\]
If $t<b$, then $t\le K_\delta$ immediately, since
$K_\delta\ge b$. If $t\ge b$, squaring the nonnegative sides gives
\[
 (1-\sigma^2)t^2-2(b+\sigma^2\delta)t
             +b^2-\sigma^2\delta^2-a^2\le0.
\]
The upper root of this upward quadratic is \eqref{eq:Kdelta}, because
\[
 (b+\sigma^2\delta)^2
 -(1-\sigma^2)(b^2-\sigma^2\delta^2-a^2)
 =\sigma^2(b+\delta)^2+(1-\sigma^2)a^2.
\]
Therefore $t\le K_\delta$. Substituting $\delta=0$ proves the final
assertions.
\end{proof}

\begin{corollary}[Entire fibers on a covered region]\label{cor:coveredfibers}
Let $\Omega\subset\dom C$ be a region covered within distance $\delta$
by the sample parameters. If every point of a nonempty fiber $F(x)$
has parameter in $\Omega$, then
$d_H(F(x),\{A_m(x)\})\le K_\delta/\lambda$. If every point of a
nonempty inverse fiber $F^{-1}(v)$ has parameter in $\Omega$, then
$d_H(F^{-1}(v),\{A_m^{-1}(v)\})\le K_\delta$.
\end{corollary}
\begin{proof}
Apply Theorem~\ref{thm:covered} to every graph point in the indicated
fiber and use the singleton formula for Hausdorff distance.
\end{proof}

The condition is coverage of $p=x+\lambda v$, not coverage of $x$
alone. For an explicit sufficient parameter radius, suppose $C(0)$ is
known and the complete Cayley map is defined. If either
$\norm{x}\le B$ or $\lambda\norm v\le B$, then
\begin{equation}\label{eq:parameterradius}
 \norm p\le
 T_B:=\max\{4B+2\norm{C(0)},(2L)^{1/(1-\gamma)}\}.
\end{equation}
Indeed, from either $p=2x-C(p)$ or $p=2\lambda v+C(p)$,
$t=\norm p\le2B+\norm{C(0)}+Lt^\gamma$. If $t>T_B$, then
$2B+\norm{C(0)}<t/2$ and $Lt^\gamma<t/2$, a contradiction.
Thus a parameter $\delta$-net of $\overline B(0,T_B)$ certifies all
fibers in either specified original-coordinate working region.
Access to such parameter samples is an additional sampling assumption,
not a consequence of having one arbitrarily chosen value at each sampled
original input.

\subsection{Finite-precision evaluation}

\begin{proposition}[Computable evaluation certificates]\label{prop:evaluation}
For a feasible $\widehat\theta\in\Delta_m$, let
$g_q=2Q\widehat\theta+d(q)$ and define its Frank--Wolfe gap
\[
 G=\ip{g_q}{\widehat\theta}-\min_i(g_q)_i.
\]
Then $G\ge0$ and
\begin{equation}\label{eq:qpgap}
 \norm{V\widehat\theta-N_m(q)}\le\sqrt G.
\end{equation}
More generally, if $\widehat N(q)$ approximates $N_m(q)$ with certified
error at most $e$, then any candidate parameter $q$ gives
\begin{align}
 \left\|A_m(x)-\frac{q-x}{\lambda}\right\|
 &\le\frac{\norm{q-2x+\widehat N(q)}+e}
              {\lambda(1-\sigma)},\label{eq:forwardresidual}\\
 \left\|A_m^{-1}(v)-(q-\lambda v)\right\|
 &\le\frac{\norm{q-2\lambda v-\widehat N(q)}+e}
              {1-\sigma}.\label{eq:inverseresidual}
\end{align}
\end{proposition}

\begin{proof}
The minimum component of $g_q$ is its minimum over the simplex, so the
gap is nonnegative. Convexity gives
$\Phi_q(\widehat\theta)-\Phi_q(\theta(q))\le G$. The exact quadratic
expansion, together with the first-order inequality at $\theta(q)$,
gives
\[
 \Phi_q(\widehat\theta)-\Phi_q(\theta(q))
 \ge(\widehat\theta-\theta(q))^TQ(\widehat\theta-\theta(q))
 \ge\norm{V\widehat\theta-N_m(q)}^2,
\]
proving \eqref{eq:qpgap}.

If $T$ is a contraction of modulus $\sigma$ with fixed point $q_*$,
then
$\norm{q-q_*}\le\norm{q-T(q)}+\sigma\norm{q-q_*}$, hence
$\norm{q-q_*}\le\norm{q-T(q)}/(1-\sigma)$.
Use $T(q)=2x-N_m(q)$ or $T(q)=2\lambda v+N_m(q)$, bound the exact
residual by its approximate residual plus $e$, and use
$A_m(x)=(q_*-x)/\lambda$ or $A_m^{-1}(v)=q_*-\lambda v$.
This gives \eqref{eq:forwardresidual}--\eqref{eq:inverseresidual}.
\end{proof}

Exact evaluations allow ordinary geometric fixed-point iteration. The
certificates above also allow any positive target accuracy using
sufficiently accurate quadratic-program solutions. Numerical errors
must be added to the geometric bounds in Theorem~\ref{thm:covered}.
Floating-point gaps need validated tolerances or outward rounding if
they are used for computer-assisted rigorous numerical certification.
No dimension-independent sample or running-time complexity is asserted;
covering a region may require many samples in high dimension.

\begin{corollary}[Bias, observation noise, and evaluation error]
\label{cor:totalerror}
Let $v$ be a true output and $\widetilde v$ an observation with
$\norm{\widetilde v-v}\le\eta$. Suppose all points of
$F^{-1}(v)$ satisfy the coverage condition of
Corollary~\ref{cor:coveredfibers}. If an evaluated proxy
$\widehat x$ satisfies
$\norm{\widehat x-A_m^{-1}(\widetilde v)}\le e$, then every
$x\in F^{-1}(v)$ satisfies
\begin{equation}\label{eq:totalerror}
 \norm{x-\widehat x}
 \le K_\delta+
       \frac{\lambda(1+\sigma)}{1-\sigma}\eta+e.
\end{equation}
\end{corollary}
\begin{proof}
Insert $A_m^{-1}(v)$ and $A_m^{-1}(\widetilde v)$ between $x$ and
$\widehat x$. Bound the three resulting distances by the coverage
certificate, the inverse Lipschitz constant in
\eqref{eq:conditioning}, and $e$, respectively.
\end{proof}

The first term is model/coverage bias, the second is noise amplification,
and the third is evaluation error. The first does not generally vanish
when observation and computation errors vanish.

\subsection{Why a finite total transcript cannot certify the whole space}

\begin{proposition}[Finite-information obstruction]\label{prop:information}
Let $n\ge1$. Consider deterministic procedures which, given point
evaluations of an unknown globally $L$-H\"older map
$C:\R^n\to\R^n$ of exponent $\gamma$, make finitely many possibly
adaptive queries and output a single-valued map $A:\R^n\to\R^n$
with no further access to the oracle. No such procedure can guarantee
a finite uniform forward-fiber
error for the corresponding relation on all of $\R^n$ for every such
$C$. In particular it cannot guarantee the sharp global radius from
finite total information alone.
\end{proposition}

\begin{proof}
Run a proposed procedure on $C_0=0$ and let $S$ contain the finitely
many queried points together with $0$. Fix a unit vector $e$ and define
$C_1(p)=L\dist(p,S)^\gamma e$. Distance to $S$ is $1$-Lipschitz, and
$\abs{s^\gamma-t^\gamma}\le\abs{s-t}^\gamma$ for nonnegative $s,t$.
Therefore $C_1$ is $L$-H\"older. It vanishes at every queried point.
By induction over the adaptive queries, the deterministic procedure
receives the same transcript on $C_1$ as on $C_0$ and outputs the same
map $A$. Both full-domain maps define graph-maximal RL relations.

The relation for $C_0$ is $F_0(x)=\{x/\lambda\}$. If the output has
no finite uniform error for $F_0$, the procedure already fails. Otherwise
there is $B<\infty$ with
$\norm{A(x)-x/\lambda}\le B$ for every $x$.
For a graph point of the relation associated with $C_1$, put
$c=C_1(p)$, $x=(p+c)/2$, and $v=(p-c)/(2\lambda)$. Then
\[
 \norm{v-A(x)}
 \ge\norm{v-x/\lambda}-\norm{A(x)-x/\lambda}
 \ge\norm c/\lambda-B.
\]
Since $S$ is finite, $\dist(p,S)$ is unbounded on $\R^n$. The last
expression is unbounded, so $A$ has no finite uniform forward error for
this second relation. This contradicts the proposed guarantee.
\end{proof}

This statement concerns finite total information, deterministic
procedures, an unrestricted H\"older class, and an unbounded domain.
It does not cover probabilistic guarantees, continued oracle access,
known analytic maps, or approximation on a specified compact region.

\section{A deadband inverse problem: what the certificate changes}
\label{sec:deadband}

Let $U$ be an orthogonal $2\times2$ matrix, let $a>0$, and put
\begin{equation}\label{eq:deadband}
 K=U[-a,a]^2,\qquad F(x)=x-P_Kx,\qquad \lambda=1.
\end{equation}
This is a vector deadband: inputs inside $K$ all produce output zero.
It is already a monotone single-valued map, and
\begin{equation}\label{eq:exactEB}
 F^{-1}(0)=K,\qquad \dist(x,K)=\norm{F(x)}.
\end{equation}
The example is not intended to demonstrate that monotone-operator
theory fails to apply. It distinguishes distance to a fixed solution
set from stable point reconstruction as the target changes.

\begin{proposition}[Exact graph coordinates and unavoidable inverse jumps]
\label{prop:deadbandjump}
The full Cayley map of \eqref{eq:deadband} is $C=P_K$. For
$\gamma=1/2$, its least H\"older constant is
$L=\sqrt D$, where $D=\diam K=2a\sqrt2$, and hence $\rad=D$.
Let $n_j=Ue_j$ and choose $0<4\eps<a$. For the two outputs
$w_\pm=\pm\eps n_1$,
\begin{equation}\label{eq:deadbandfibers}
 F^{-1}(w_\pm)=\{\pm(a+\eps)n_1+tn_2:\abs t\le a\}.
\end{equation}
Every pair of exact reconstructions $x_\pm\in F^{-1}(w_\pm)$ satisfies
\begin{equation}\label{eq:exactgain}
 \frac{\norm{x_+-x_-}}{\norm{w_+-w_-}}
 \ge\frac a\eps+1.
\end{equation}
\end{proposition}

\begin{proof}
If $q=P_Kx$, then the projection characterization says
$\ip{x-q}{y-q}\le0$ for every $y\in K$. Thus
$p=2x-q$ has $P_Kp=q$, and
$x-F(x)=q=P_K(x+F(x))$. Conversely, for any $p$ set
$q=P_Kp$ and $x=(p+q)/2$. The same projection characterization,
scaled by $1/2$, gives $P_Kx=q$ and $F(x)=(p-q)/2$.
This proves the full parametrization, hence graph maximality.
Nonexpansiveness and the image diameter give
\[
 \norm{P_Kp-P_Kq}\le\min\{\norm{p-q},D\}
 \le\sqrt D\,\norm{p-q}^{1/2}.
\]
Two diameter endpoints of $K$, fixed by $P_K$, force
$L\ge\sqrt D$, so the constant is least possible.

In rotated coordinates, each component of $F$ is
$f(s)=s-\max\{-a,\min\{s,a\}\}$. For an output $\eps>0$ its
unique scalar inverse is $a+\eps$; for $-\eps$ it is $-a-\eps$;
for zero its inverse is $[-a,a]$. This yields
\eqref{eq:deadbandfibers}. The difference in the $n_1$ coordinate
of any two chosen inverse points is $2a+2\eps$. Their Euclidean
distance is at least that amount, while
$\norm{w_+-w_-}=2\eps$, proving \eqref{eq:exactgain}.
\end{proof}

For a transparent analytic proxy, take the contraction
$N=\tfrac12 P_K$ and define $A$ by \eqref{eq:shadowparam}. Its
Lipschitz and inverse Lipschitz constants are at most $3$, and its
strong-monotonicity modulus is at least $1/3$.
This is a model-specific contraction; its choice $\sigma=1/2$ is not
the radius-minimizing parameter $\sqrt\gamma=1/\sqrt2$ of the
general theorem.

\begin{proposition}[Stable analytic proxy with full-fiber control]
\label{prop:deadbandproxy}
For the above $A$, in the rotated coordinates each component of
$A^{-1}$ is
\begin{equation}\label{eq:deadbandproxyformula}
 \alpha(w)=
 \begin{cases}
 3w,& |w|\le a/4,\\
 w+\tfrac a2\operatorname{sign}(w),& |w|\ge a/4.
 \end{cases}
\end{equation}
The formulas agree at the endpoints. Every complete inverse fiber obeys
\begin{equation}\label{eq:deadbandglobaltube}
 \sup_{x\in F^{-1}(v)}\norm{x-A^{-1}(v)}\le a\sqrt2
 \qquad(v\in\R^2).
\end{equation}
For $w_\pm$ from Proposition~\ref{prop:deadbandjump},
\begin{align}
 A^{-1}(w_\pm)&=\pm3\eps n_1,\label{eq:proxygain1}\\
 \frac{\norm{A^{-1}(w_+)-A^{-1}(w_-)}}
      {\norm{w_+-w_-}}&=3,\label{eq:proxygain2}\\
 \sup_{x\in F^{-1}(w_\pm)}\norm{x-A^{-1}(w_\pm)}
 &=\sqrt{(a-2\eps)^2+a^2}.\label{eq:proxyerror}
\end{align}
\end{proposition}

\begin{proof}
Writing $\pi(p)=\max\{-a,\min\{p,a\}\}$, scalar inverse evaluation
requires $p=2w+\tfrac12\pi(p)$ and returns $p-w$.
If $|w|\le a/4$, then $p=4w\in[-a,a]$ solves this equation and
returns $3w$. If $w\ge a/4$, then $p=2w+a/2\ge a$ solves it and
returns $w+a/2$; the negative case is symmetric. The contraction
ensures these are the unique solutions, proving
\eqref{eq:deadbandproxyformula}.

For $w>0$, the scalar true inverse is $a+w$. If $0<w\le a/4$,
its difference from the proxy is $a-2w\le a$; if $w\ge a/4$, its
difference is $a/2$. For $w<0$ the absolute differences are the
same. At $w=0$, the true inverse interval is $[-a,a]$ and the proxy
is zero, so the worst difference is $a$. Applying this bound to
each rotated coordinate and using orthogonality gives
\eqref{eq:deadbandglobaltube}. Since $4\eps<a$, the first branch
of \eqref{eq:deadbandproxyformula} applies to both targets, proving
\eqref{eq:proxygain1}--\eqref{eq:proxygain2}. Substitution into
\eqref{eq:deadbandfibers} gives the displacement
$\pm(a-2\eps)n_1+tn_2$; maximizing its norm over $|t|\le a$ proves
\eqref{eq:proxyerror}.
\end{proof}

For $a=1/8$ and $\eps=10^{-4}$, the output displacement is $0.0002$.
Every exact inverse pair moves by at least $0.2502$, a gain of at
least $1251$, whereas the analytic proxy moves by $0.0006$, a gain
of $3$. The proxy generally is not an exact inverse point. These
numbers compare two different requirements---exact reconstruction and
certified approximate reconstruction---and make their tradeoff explicit.

The known analytic model permits other simple choices. In particular,
the identity map already gives the global forward and inverse radius
$a\sqrt2$, because $x-v=P_Kx\in K$ and
$\sup_{k\in K}\norm k=a\sqrt2$. Thus this example does not establish
that the shadow construction dominates bespoke regularization. Its
purpose is to exhibit what a common stable proxy and a full-fiber
certificate mean, and what they do not mean. For finite graph data,
Section~\ref{sec:finite} supplies a model-class construction and adds
the explicit coverage and evaluation terms that an exact analytic
calculation does not require.

\section{Discussion and scope}\label{sec:discussion}

The simultaneous shadow theorem and the realization theorem express
two complementary conclusions. The entire relation admits one
quantitatively controlled, strongly monotone, invertible approximation.
Nevertheless, an individual finite-dimensional fiber may be any
nonempty compact set of the permitted diameter. The approximation
does not turn original fibers into convex sets, select an exact branch,
or eliminate nonidentifiability. Its radius is a uniform worst-case
guarantee, not a claim that every relation requires the full radius.

The global assumption is restrictive. To see its large-scale meaning,
fix a graph anchor and write $d=(x-x_0)+\lambda(v-v_0)$ and
$e=(x-x_0)-\lambda(v-v_0)$. Along any graph sequence unbounded in
$H\times H$, $\norm d$ must be unbounded, since
$\norm e\le L\norm d^\gamma$ bounds both original coordinate
differences whenever $d$ is bounded. Along a sequence with
$\norm d\to\infty$, one has $\norm e/\norm d\to0$.
Since $x-x_0=(d+e)/2$, this implies, uniformly for large graph points,
\[
 v-v_0=\lambda^{-1}(x-x_0)+O(\norm{x-x_0}^\gamma).
\]
In particular, the class prescribes a common asymptotic linear slope;
it is not a container for every locally regular nonmonotone model.

The sharp constant established here is uniform over dimensions.
Determining optimal constants in each fixed dimension, or under
additional geometric assumptions on fibers, remains a distinct
question. Likewise, a specified anchor changes the approximation
problem and its optimal factor. Additional interpolation constraints
or stronger computational requirements may impose different tradeoffs.

The finite-data theorem distinguishes a mathematical existence result
from a computable surrogate. It supplies one coherent global map from
fixed observations and certifies unobserved graph points only where
the specified coverage assumption holds. Proposition~\ref{prop:information}
shows why this locality cannot be discarded in the stated oracle
model. Efficient adaptive coverage, structured sampling, and
dimension-dependent guarantees are natural computational questions,
but are not resolved by the existence of a finite quadratic program.

\begin{thebibliography}{99}

\bibitem{AzagraLeGruyerMudarra}
D. Azagra, E. Le Gruyer, and C. Mudarra,
\emph{Kirszbraun's theorem via an explicit formula},
Canadian Mathematical Bulletin \textbf{64}(1) (2021), 142--153.
\url{https://doi.org/10.4153/S0008439520000314}.

\bibitem{BauschkeWang}
H. H. Bauschke and X. Wang,
\emph{Firmly nonexpansive and Kirszbraun--Valentine extensions: a
constructive approach via monotone operator theory},
in Nonlinear Analysis and Optimization I, Contemporary Mathematics
\textbf{513}, American Mathematical Society, 2010, 55--64.
\url{https://doi.org/10.1090/conm/513/10075}.

\bibitem{BenevieriFuriPeraSpadini}
P. Benevieri, M. Furi, M. P. Pera, and M. Spadini,
\emph{An introduction to topological degree in Euclidean spaces},
arXiv:2304.06463, 2023. \url{https://arxiv.org/abs/2304.06463}.

\bibitem{BurachikIusemSvaiter}
R. S. Burachik, A. N. Iusem, and B. F. Svaiter,
\emph{Enlargement of monotone operators with applications to variational
inequalities}, Set-Valued Analysis \textbf{5}(2) (1997), 159--180.
\url{https://doi.org/10.1023/A:1008615624787}.

\bibitem{Kirszbraun}
M. D. Kirszbraun,
\emph{\"Uber die zusammenziehende und Lipschitzsche Transformationen},
Fundamenta Mathematicae \textbf{22}(1) (1934), 77--108.
\url{https://doi.org/10.4064/fm-22-1-77-108}.

\bibitem{LevyRice}
R. Levy and M. D. Rice,
\emph{The approximation and extension of uniformly continuous Banach
space valued mappings}, Commentationes Mathematicae Universitatis
Carolinae \textbf{24}(2) (1983), 251--265.
\url{https://dml.cz/handle/10338.dmlcz/106224}.

\bibitem{LukeTam}
D. R. Luke and M. K. Tam,
\emph{Generalized monotonicity and the proximal point algorithm},
Mathematics of Operations Research, Articles in Advance,
published online 21 November 2025.
\url{https://doi.org/10.1287/moor.2025.0863}.

\bibitem{Minty1970}
G. J. Minty,
\emph{On the extension of Lipschitz, Lipschitz--H\"older continuous,
and monotone functions}, Bulletin of the American Mathematical Society
\textbf{76}(2) (1970), 334--339.
\url{https://doi.org/10.1090/S0002-9904-1970-12466-1}.

\bibitem{RubbensHendrickxTaylor}
A. Rubbens, J. M. Hendrickx, and A. Taylor,
\emph{A constructive approach to strengthen algebraic descriptions of
function and operator classes},
Journal of Nonsmooth Analysis and Optimization \textbf{6} (2026),
article 15673.
\url{https://doi.org/10.46298/jnsao-2026-15673}.

\bibitem{Schoenberg}
I. J. Schoenberg,
\emph{Metric spaces and positive definite functions},
Transactions of the American Mathematical Society \textbf{44}(3) (1938),
522--536.
\url{https://doi.org/10.1090/S0002-9947-1938-1501980-0}.

\bibitem{Valentine}
F. A. Valentine,
\emph{A Lipschitz condition preserving extension for a vector function},
American Journal of Mathematics \textbf{67}(1) (1945), 83--93.
\url{https://doi.org/10.2307/2371917}.

\bibitem{ZaichykEtAl}
H. Zaichyk, A. Biess, A. Kontorovich, and Y. Makarychev,
\emph{Efficient Kirszbraun extension with applications to regression},
Mathematical Programming \textbf{207} (2024), 617--642.
\url{https://doi.org/10.1007/s10107-023-02023-6}.

\end{thebibliography}
\end{document}
